\section{模型假设}
\label{sec:assumptions}

\begin{enumerate}
  \item \textbf{指标方向的代理假设}。16 项入模质量指标中，无字母词比例、最高频二元/三元词组占比按负向处理，作为噪声与重复性的代理。数学、表格、代码类文本中的合理重复可能受此惩罚，该假设的影响已通过删除负向指标的对照方案量化（第~\ref{sec:q1-quality}~节，秩相关 $0.967$）。
  \item \textbf{冻结标尺假设}。归一化边界（$A_1$ 拟合集 $1\%/99\%$ 分位数）与 CRITIC 权重仅由 $A_1$ 拟合集估计，固定用于留出集与扩展集，不随新数据重估，保证全部样本在同一标尺上可比。
  \item \textbf{完全可补偿与有限补偿的分工假设}。TOPSIS 总分对 16 指标完全可补偿，用于建立可比基线；冲突消解模型对五维结构引入有界补偿（$\lambda\le1$），用于识别与下调“优缺点相抵”型样本。两者是递进设计，原评分始终作为后续三问的输入。
  \item \textbf{五维分组的适用性假设}。22 项质量指标按语义压缩为五维（内容价值、表达质量、整洁度、无广告程度、非重复程度），组内继承 CRITIC 权重；DSIR 三项为指定目标域亲和度，不计入五维。该分组是建模选择，冲突率水平对分组方案敏感（第~\ref{sec:q1-conflict}~节已声明）。
  \item \textbf{尺度分层假设}。配比—损失关系在同尺度内建模，跨尺度只输出相对排序类结论；不做只含配比的模型的跨尺度绝对损失预测（由重复设计的配对分析支持）。
  \item \textbf{半合成数据的证据分级假设}。B2、B6--B8 等半合成数据仅用于补充分析与敏感性讨论，不作为直接实验观测；C6 桥接按 Loss\_Comparability 分层使用，高可比与中可比样本不混同推断。
  \item \textbf{外生变量假设}。问题三中上下文长度 $S$ 由模型架构与任务需求外生给定，其可行取值依据 C7，通过注意力开销项影响预算，不作为内点寻优变量。
\end{enumerate}
